% Part: lambda-calculus
% Chapter: introduction
% Section: syntax

\documentclass[../../../include/open-logic-section]{subfiles}

\begin{document}

\olfileid{lam}{int}{syn}
\olsection{లాంబ్డా కలనశాస్త్ర వాక్యనిర్మాణం}

$x$, $y$, $z$,~\dots అనే చరాల క్రమంతోనూ, $a$, $b$, $c$,~\dots
అనే కొన్ని స్థిరసంకేతాలతోనూ ప్రారంభిస్తాం. పదాల సమితిని కింది విధంగా
ఆగమనాత్మకంగా నిర్వచిస్తాం:
\begin{enumerate}
\item ప్రతి చరం ఒక పదం.
\item ప్రతి స్థిరసంకేతం ఒక పదం.
\item $M$, $N$ పదాలైతే, $(MN)$ కూడా ఒక పదం.
\item $M$ ఒక పదమూ, $x$ ఒక చరమూ అయితే, $(\lambd[x][M])$ ఒక పదం.
\end{enumerate}
$(MN)$ రూపంలోని పదాలను \emph{ప్రయోగాలు} అనీ,
$(\lambd[x][M])$ రూపంలోని వాటిని \emph{అమూర్తీకరణలు} అనీ అంటాం.

స్థిరసంకేతాలే లేని వ్యవస్థను \emph{శుద్ధ} లాంబ్డా కలనశాస్త్రం అంటారు.
మనం ప్రధానంగా శుద్ధ $\lambd$-కలనశాస్త్రంలో పనిచేస్తాం; అందువల్ల చిన్న
అక్షరాలన్నీ చరాలను సూచిస్తాయి. $\lambd$-కలనశాస్త్ర పదాలను సూచించడానికి
పెద్ద అక్షరాలను ($M$, $N$ మొదలైనవి) వాడతాం.

మనం కొన్ని సంకేతన సంప్రదాయాలను పాటిస్తాం:

\begin{conv}
\begin{enumerate}
\item కుండలీకరణలను విడిచిపెట్టినప్పుడు, ప్రయోగం ఎడమ నుంచి కుడికి
  జరుగుతుంది. ఉదాహరణకు, $M$, $N$, $P$, $Q$ పదాలైతే,
  $MNPQ$ అనేది $(((MN)P)Q)$కు సంక్షిప్త రూపం.
\item మళ్లీ, కుండలీకరణలను విడిచిపెట్టినప్పుడు, లాంబ్డా అమూర్తీకరణకు
  సాధ్యమైనంత విస్తృత వ్యాప్తిని ఇవ్వాలి. ఉదాహరణకు,
  $\lambd[x][MNP]$ను $(\lambd[x][((MN)P)])$గా చదువుతాం.
\item అనేక చరాలను అమూర్తీకరించడానికి ఒకే లాంబ్డాను వాడవచ్చు.
  ఉదాహరణకు, $\lambd[xyz][M]$ అనేది
  $\lambd[x][\lambd[y][\lambd[z][M]]]$కు సంక్షిప్త రూపం.
\end{enumerate}
\end{conv}

ఉదాహరణకు,
\[
\lambd[xy][xxyx \lambd[z][xz]]
\]
అనేది
\[
\lambd[x][\lambd[y][((((xx)y)x)(\lambd[z][(xz)]))]].
\]
కు సంక్షిప్త రూపం. ఈ సంప్రదాయాలను గుర్తుంచుకోవాలి. మొదట ఇవి మీకు
చికాకు కలిగిస్తాయి; కానీ క్రమంగా అలవాటవుతాయి. కొంతకాలానికి, గజిబిజి
కుండలీకరణలతో వ్యవహరించాల్సి రావడం కంటే ఇవే తక్కువ చికాకు కలిగిస్తాయి.

బద్ధ చరాల పేర్లలో మాత్రమే భేదమున్న రెండు పదాలను $\alpha$-తుల్యమైనవి
అంటాం; ఉదాహరణకు, $\lambd[x][x]$, $\lambd[y][y]$. వీటిని ``ఒకే''
పదంగా భావించడం అనుకూలం; మరో మాటలో, $M$, $N$ ఒకటే అని చెప్పినప్పుడు,
``బద్ధ చరాల పేరుమార్పుల వరకు'' ఒకటే అని కూడా ఉద్దేశిస్తాం. ఒక
$\lambd$ వ్యాప్తిలో ఉన్న చరాలను ``బద్ధమైనవి'' అనీ, మిగిలినవాటిని
``స్వేచ్ఛావి'' అనీ అంటాం. ముందున్న ఉదాహరణలో స్వేచ్ఛా చరాలు లేవు;
కానీ
\[
(\lambd[z][yz])x
\]
లో $y$, $x$ స్వేచ్ఛా చరాలు; $z$ బద్ధ చరం.

\end{document}
